\documentclass[14pt, oneside]{article} 
\usepackage{geometry}
\geometry{letterpaper} 
\usepackage{graphicx}
	
\usepackage{amssymb}
\usepackage{amsmath}
\usepackage{parskip}
\usepackage{color}
\usepackage{hyperref}

\graphicspath{{/Users/telliott/Github-math/figures/}}
% \begin{center} \includegraphics [scale=0.4] {gauss3.png} \end{center}

\title{Feuerbach inversively}
\date{}

\begin{document}
\maketitle
\Large

%[my-super-duper-separator]

\begin{center} \includegraphics [scale=0.55] {inversion16.png} \end{center}  

In this chapter we are going to look at Feuerbach's theorem involving the nine point circle.  We will use inversive geometry to do this.  We have already introduced the basic principles, here we must expand on a few areas before we attack the main problem.

We recall the fundamental definition of inversive geometry.
\begin{center} \includegraphics [scale=0.40] {inversion1.png} \end{center}
\[ OP \cdot OP' = k^2 \ \ \ \ \ \ \ \  \frac{OP}{k} = \frac{k}{OP'} \]
The image of a point $P$ under inversion is called $P'$, and it lies on the same ray as $OP$ such that it satisfies the above equation.

\subsection*{Invariants in inversive geometry}

Either a given point is on a circle (or line) or it is not.  If it is on the locus, then it must also lie on the image of the locus under inversion.

This seems quite obvious, yet subtle.  We will not prove it here.

Then if a circle or a line cuts a circle at two points, after inversion the resulting figures will still have two points of intersection.

And if a line is tangent to a circle or if two circles are tangent, or a line cuts a line, after inversion the resulting figures will still be tangent (in the first two cases) and in all cases, share one point of intersection.

The next figure shows a circle $K$ not through $O$ and its inverse $K'$.  We notice that lines tangent to the original circle $K$ are also tangent to the inverse circle $K'$

However, although the center of $K$ and the center of $K'$ are on the same line through $O$, they are \emph{not} inverses.

\begin{center} \includegraphics [scale=0.35] {inversion10.png} \end{center} 

We'll see why in just a bit.

\subsection*{Inversion preserves angles}.

As a special case (and the only one we will use for the Feuerbach theorem, below), if two circles are orthogonal, then their inverses are also orthogonal.

Two circles are orthogonal if, at the points where they cross, the two tangents are perpendicular.  Equivalently, the two radii to that point must be perpendicular.

Then, considering a circle orthogonal to $\omega$, the claim is that inversion in $\omega$ leaves such a circle invariant, or fixed.

\begin{center} \includegraphics [scale=0.45] {FB4.png} \end{center}

\emph{Proof}.

Let the inversion circle $\omega$ have its center at $O$ and have radius of length $k$.

Another circle $K$ with radius of length $R$ is positioned so that at the points where the two circles cross, the tangents (and radii) are perpendicular.

The first thing we notice is that the points $S$ and $T$ lie on $\omega$, so they are fixed, their inverses are the same as the original point.

Next consider a pair of points lying $P$ and $P'$ on a diameter of $K$, $P'$ on a distance $x$ from the inversion center, while $P$ lies a distance $2R+x$ from the inversion center.  The product of the distances is
\[ (2R + x)x = 2Rx + x^2 \]

But we also have from the geometry that
\[ k^2 + R^2 = (R+x)^2 \]
\[ k^2 = 2Rx + x^2 \]

The product of the two distances is just $k^2$, so the inverse of $P$ is $P'$ and vice-versa.  We now know that four points of the inverted circle are also on the original circle.  The result follows.

Even more generally, we may consider an arbitrary secant $QQ'$.  Recall from the tangent-secant theorem that $OQ \cdot OQ' = OT^2 = k^2$.

Any pair of points on any secant of circle $K$ are their own inverses, so they both lie on the inverted circle.  Thus, the circle is fixed under inversion.

$\square$

A relatively simple proof for circles that are not necessarily orthogonal comes from the web.  I have lost the author's name.

We start with the special case where one circle goes through $O$:

\begin{center} \includegraphics [scale=0.30] {inversion11.png} \end{center}

The line $m$ through $O$ transforms into itself.  The circle is also through $O$, and is the inverse of the line $l$.  The angle between $l$ and $m$ is easily shown to be equal to the angle between $m$ and the image of the line $l$, namely the circle, at $A'$.

The shaded gray angle at $A$ is complementary to $\angle AOP$.  But $\angle AOP = \angle QOA' = \angle QA'O$, and $\angle QA'O$ is complementary to the gray shaded angle at $A'$.  $\square$ 

In the general case, suppose the circles on centers $B$ and $B'$ are inverses through $I_{\omega}$.  Let $A$ and $A'$ be inverse points as usual.

\begin{center} \includegraphics [scale=0.30] {inversion12.png} \end{center}

This does \emph{not} mean that $I_{\omega} (B') = B$!

However, the green angles at $A'$ and $A_1$ \emph{are} equal by similar triangles.  This requires a bit of work.  We simplify the diagram and change the notation as far as the center of the circle, and also introduce tangent points $T$.

\begin{center} \includegraphics [scale=0.25] {inversion14.png} \end{center}

By the secant-tangent theorem, $OA_1 \cdot OA' = OT_1^2$ and this is true wherever $A'$ and $A_1$ lie.

By the definition of inversion $OA \cdot OA' = k^2$.  So

\[ \frac{OA_1}{OA} = \frac{OT_1^2}{k^2} \]
which is a constant.  Call it $c$.  So then $OA_1 = c \cdot OA$ for every pair of points $A_1$ and $A$ (\emph{at the second point of intersection}) lying on the same line through circles $Q$ and $Q_1$.

\begin{center} \includegraphics [scale=0.25] {inversion14.png} \end{center}

Since $\triangle OQT \sim \triangle OQ_1T_1$, that includes $Q$ and $Q_1$ so $OQ_1 = c \cdot OQ$. 

This relationship is called a homothety.  The entire circle on $Q_1$ is a bigger version of the one on $Q$.

It follows that $\triangle OAQ \sim \triangle OA_1 Q_1$.  We now switch back to the original notation:  $\triangle OA'B' \sim \triangle OA_1B$, which gives two of the equal angles colored green in the figure below.  The third one (at $A$) is equal since $\triangle BAA_1$ is isosceles.

\begin{center} \includegraphics [scale=0.20] {inversion13.png} \end{center}

The gray angles are both complementary to the green angles.  The second one in circle $B$ (at $A$) is equal by reflection across a line of symmetry through $B$ and $\perp m$.

$\square$

\begin{center} \includegraphics [scale=0.25] {Coxeter_5_5A.png} \end{center}

\emph{Proof}.  (Alternate).

This is section 5.5 of Coxeter.

Consider two circles intersecting at $P$ (not drawn).  The supplementary angles formed between the circles are naturally defined in terms of the tangents to the two circles at $P$.  Suppose the angles are $\theta$ and $\phi$.  Now consider just the two tangent lines $a$ and $b$, which form $\angle \theta$ in the figure above.

We perform an inversion with respect to circle $\omega$ on center $O$.  The inverse of a line not through $O$ is a circle through $O$, where the tangent to the new circle at $O$ is parallel to the original line.  So the inverse of line $a$ is circle $\alpha$ with its tangent parallel to $a$, and similarly for $b$.  We can see that the angle between these two tangents at $O$ is also equal to $\theta$, since they are parallel to the original lines $a$ and $b$.

As we said above, if $P$ lies on both circles before inversion, it lies on both circles (or lines) after inversion.  Thus the tangents at $P'$ are equal to those at $P$ and form the same angle.

We can also see that for intersecting circles, the angles at the two points of intersection are the same although the tangents have a mirror image symmetry.
\begin{center} \includegraphics [scale=0.45] {FB5a.png} \end{center}

\subsection*{excircles}

To review briefly, \emph{excircles} are so named by reference to \emph{incircles}.  The incircle to a triangle is tangent to all three sides and lies inside the triangle.  An excircle is tangent to one side and also tangent to extensions of the other two sides and it is external to the triangle.

In $\triangle ABC$, with opposite sides $a$, $b$ and $c$, the incircle on center $M$ is that circle to which all three sides are tangent.

\begin{center} \includegraphics [scale=0.3] {excircle.png} \end{center}  

Any triangle has a single incircle and three excircles, one for each side of the triangle.  The excircle on side $c$ (opposite $\angle C$) is formed by extending sides $a$ and $b$ (i.e. $BC$ and $AC$ above) and then finding the circle that is tangent to those two extensions and also tangent to side $c$.

\begin{center} \includegraphics [scale=0.15] {heron7.png} \end{center}
Let $N$ be the center of the excircle on side $c$.  $N$ lies on the bisector of the original $\angle C$.  This follows from the definition that $CD$ and $CE$ are tangent to this excircle.

Since $AD$ and $AJ$ are also tangents to the excircle from point $A$, if the angle supplementary to $\angle A$ (i.e. $\angle BAD$) is bisected, then the bisector also runs through $N$, so $N$ can be constructed by finding the intersection of the bisectors of $\angle DAB$ and $\angle C$.  The bisector of $\angle ABE$ goes through $N$ as well.

$s$ is the \emph{semi-perimeter} of the triangle, one-half the sum of the side lengths. $AG = u, BF = v, FC = w$.  As tangents from a point, the perimeter of the triangle has two copies of each, so the semi-perimeter has one:
\[ u + v + w = s \]
But $v + w = a$, so $AG = u = s - a$ and so on.

It is also easy to show that the tangents from $A$ to the excircle on side $c$ are 
\[ AD = AJ = s-b \]
In other words, $J$ and $H$ each divide side $c$ into lengths $s-a$ and $s-b$.  \emph{Proof}.  Let $AD = p$ and $BE = q$. Then
\[ p + (s-a) + (s-c) = q + (s-b) + (s-c) \]
\[ p - a = q - b \]
But we also know that $p + q = c$ so
\[ c - q - a = q - b \]
\[ 2q = - a + b + c = 2(s-a) \]
\[ q = BE = BJ = s-a \]
$\square$
\begin{center} \includegraphics [scale=0.15] {heron7.png} \end{center}
The length of the long tangents $CD = CE$ is easily shown to be just $s$.

\[ (s-a) + (s-b) + (s-c) = 3s - (a + b + c) = s \]

\subsection*{Feuerbach's Theorem}

Finally, we reach this famous theorem.

\begin{center} \includegraphics [scale=0.15] {Feuerbach.png} \end{center} 

Recall that the nine point circle, passing through the midpoints of the sides, is the circumcircle of the medial triangle, i.e., it is formed with these points as the vertices.

Feuerbach's theorem says that this circle (in red below) is tangent (internally) to the incircle and (externally) to each of the excircles of a triangle.

\begin{center} \includegraphics [scale=0.36] {FB8.png} \end{center}

We previously looked at the incircle and three excircles of a triangle.  Recall that the nine point circle also goes through the feet of the altitudes and the bisectors of that part of each altitude between the orthocenter and each vertex.

Let $\triangle ABC$ have the sides bisected at $A'$, $B'$ and $C'$.  The incircle is drawn on center $I$, and one excircle is drawn tangent to side $a = BC$, on center $O_a$.  $\triangle A'B'C'$ is also drawn.

\begin{center} \includegraphics [scale=0.35] {FB1.png} \end{center}

\textbf{I}

We know from previous work that $BX = BY = s-b$, $AY = AZ = s-a$ and $CX = CZ = s-c$.  As well, the entire tangent from $A$ to the excircle has length $s$, the semi-perimeter.

Radii are drawn from the incenter $I$ and the excenter $O_a$, perpendicular to side $BC$.  $BC$ is tangent to the two circles at these points.  It follows that
\[ IX \parallel O_a X_a \]

We also know from previous work with excircles that $BX = X_a C = s-b$.  So
\[ X X_a = a - 2(s-b) = a - (a - b + c) = b - c  \]

We also have that $XA' = X_a A'$.  That's because $A'$ is the midpoint of side $BC = a$, and we've subtracted equal lengths from both ends with $BX = X_a C$.

Therefore, we can draw the circle as shown, with $A'$ as the center and radius $A'X = A'X_a = (b-c)/2$.

We will refer to that circle as $\omega$.  It will be the inversion center for the transformation.

We also draw another line tangent to both incircle and excircle.  Label the endpoints as $B_1$ and $C_1$.

Also mark three other points:   where $A'B'$ crosses $B_1 C_1$ at $B''$, where $A'C'$ crosses $B_1 C_1$ at $C''$, and finally, the point where $BC$ crosses $B_1 C_1$, at $S$.

\begin{center} \includegraphics [scale=0.35] {FB2.png} \end{center}

\textbf{II}

The center of the inversion circle $\omega$ is at $A'$, which is also on the nine-point circle.  The inverse of a circle through the inversion center is a line.

The claim is that $B''$ is the inverse of $B'$, and $C''$ is the inverse of $C'$ so the straight line $B_1 C_1$ which goes through those two points is the inverse of the nine point circle containing points $A'$, $B'$ and $C'$.

To begin with we obtain some more lengths.  Use the angle bisector theorem to get an expression for the length of $BS$ and $CS$, the components of side $a$ produced by the bisector.  The claim is that
\[ BS = \frac{ac}{b+c} \ \ \ \ \ \ CS = \frac{ab}{b + c} \]

This looks plausible because the sum is just $a = BC$, which checks.  

Coxeter explain it thus:  \begin{quote}Since $S$ (like $I$ and $I_a$) lies on the bisector of the angle $A$, $S$ divides the segment $BC$, of length $a$, in the ratio $b:c$. \end{quote}

Thus the two parts of $a$ are the fractions $b/(b+c)$ and $c/(b+c)$.

\begin{center} \includegraphics [scale=0.35] {FB2.png} \end{center}

Working through the algebra just to confirm this, let $BS = u$ and $CS = v$.  Then $u/c = v/b$ so $u/v = c/b$.  Add one to both sides
\[ \frac{u+v}{v} = \frac{b + c}{b} \]
but $u+v = a$ so we have $v = ab/(b+c)$.  Alternatively add one to both sides of the inverse of what's above.  Then
\[ \frac{u+v}{u} = \frac{b + c}{c} \]
and the second result follows easily as well.

Another claim is that
\[ SA' = \frac{a(b-c)}{2(b+c)} \]
We get this from $BC = a$, so $BA' = a/2$ and then
\[ SA' = \frac{a}{2} - \frac{ac}{b+c} \]
Placed over a common denominator and subtracting, that checks.

The last claim, which I missed in a youtube video of this proof but is in the source (Coxeter), is that $BC_1 = B_1C = b-c$.  This is easier to see for $B_1 C$ so let us start there.  The whole of side $AC$ has length $b$. so somehow it must be the case that $AB_1$ has length $c$ and then $B_1 C$ is the difference.

\begin{center} \includegraphics [scale=0.35] {FB3.png} \end{center}

We must show that $B_1 C = BC_1 = b-c$.  There is a lot of symmetry in the figure.  Both circles have their centers on the bisector of $\angle A$.

We have that $\triangle SIX$ and $\triangle SIX_1$ have three sides equal (shared side, radius, and two tangents from a point), so they are congruent, which gives $\angle ASX = \angle ASX_1$.

It follows that $\triangle ASB$ and $\triangle ASB_1$ are equiangular and share side $AS$, so they are congruent.  Thus $AB = AB_1 = c$.

So $B_1C = b - c$, by subtraction, and 
\[ B_1Z = b - (s-a) - (b-c) = a + c - s = s - b \]
One can also get this from $SB = SB_1$ (congruent triangles, above) and $SX = SX_1$, so by subtraction, $BX = B_1X_1 = B_1Z$.

The whole $\angle ASC = \angle ASC_1$ because the parts are equal (from congruent triangles and then vertical angles).  It follows that $\triangle ASC \cong \triangle ASC_1$.  So $AC = AC_1$ and $AB = AB_1$, and then finally $B_1 C = BC_1$, also by subtraction.

We do not need it but we can also show that $BC = B_1C_1$.  $\triangle ABC \cong \triangle AB_1C_1$.

\textbf{III}

\begin{center} \includegraphics [scale=0.35] {FB2.png} \end{center}
Next, we must find two pairs of similar triangles.  These are
\[ \triangle SA'B'' \sim \triangle SBC_1 \ \ \ \ \ \ \triangle SA'C'' \sim \triangle SCB_1 \]

These are easy because we have parallel lines and either vertical angles in the first case, or shared $\angle CSB_1$ in the second.

We have 
\[ \frac{SA'}{A'B''} = \frac{SB}{BC_1} \]
\[ A'B'' = SA' \cdot BC_1 \cdot \frac{1}{SB} \]
Plugging in the results from above
\[ = \frac{a(b-c)}{2(b+c)} \cdot (b-c) \cdot \frac{b+c}{ac} \]
and since $A'B' = c/2$:
\[ A'B' \cdot A'B'' = (\frac{b-c}{2})^2 \]

For the second triangle we obtain:
\[ \frac{SA'}{A'C''} = \frac{SC}{B_1C} \]
\[ A'C'' = SA' \cdot B_1C \cdot \frac{1}{SC} \]
And again, plugging in:
\[ = \frac{a(b-c)}{2(b+c)} \cdot (b-c) \cdot \frac{b+c}{ab} \]
and since $A'C' = b/2$:
\[ A'C' \cdot A'C'' = (\frac{b-c}{2})^2 \]

But this is (the same for both), namely $r^2$.

\begin{center} \includegraphics [scale=0.35] {FB2.png} \end{center}

So $B''$ and $C''$ are the inverses of $B'$ and $C'$ respectively, with respect to $\omega$.

\textbf{IV}

We have an inversion circle centered on $A'$ with diameter $XX_a$ that we have labeled $\omega$.

We have shown that the inverse of the nine-point circle is the line $B_1 C_1$, which passes through $BC$ at $S$ and is tangent to both the incircle and to the excircle (these points are not labeled).

Since the inverses of the incircle and excircle are orthogonal with respect to $\omega$, they are fixed (they invert into themselves).

So then, finally, in the inverted system we have line $B_1 C_1$ and the incircle and excircle, with the line tangent to both of these circles.

Points where two curves cross or touch are still points where the inverted figures cross or touch, after inversion.

Since $B_1 C_1$ is tangent to both circles, its inverse is still tangent to both circles.  The inverse of $B_1 C_1$ is the nine point circle, so that circle is tangent to both the other circles, i.e. the incircle and the excircle, and to the other two excircles, by extension.

\begin{center} \includegraphics [scale=0.36] {FB8.png} \end{center}

\subsection*{other ideas}

Finally, a note about other approaches.  One can show that two circles are tangent externally (such as the excircle and nine-point circle), by showing that the distance between the two centers is equal to the sum of the radii.

For internal tangency, the distance between centers must be the difference of the two radii.

Thus, if one can calculate \emph{where} the centers lie, then this requirement can be checked.  A judicious choice of coordinate system will help, and you can find such proofs on the internet.  However, I think the inversion proof presented in detail here is a lot more fun, and quite simple, once we understand what inversion is all about.


\end{document}